\QHeading{第1题　脱式计算。}
\par\noindent


\CalcBlock[方法一]{1}{\begin{aligned}[t]&\phantom{={}}\color{black}24\times[270\div(8+1)]-30\\&\color{answerred}=24\times[270\div9]-30\\&\color{answerred}=24\times30-30\\&\color{answerred}=720-30\\&\color{answerred}=690\end{aligned}}
\par
\MethodBlock{方法二：先处理括号或乘除，再提取相同因数}{用乘法分配律反过来计算。}{\begin{aligned}[t]&\phantom{={}}24\times [270\div (8+1)]-30\\
&=24\times [270\div 9]-30\\
&=24\times 30-30\\
&=30\times (24-1)\\
&=30\times 23\\
&=690\end{aligned}}


\CalcBlock{2}{\begin{aligned}[t]&\phantom{={}}\color{black}[600-(120+80)]\div25\times4\\&\color{answerred}=[600-200]\div25\times4\\&\color{answerred}=400\div25\times4\\&\color{answerred}=16\times4\\&\color{answerred}=64\end{aligned}}
\CalcBlock{3}{\begin{aligned}[t]&\phantom{={}}\color{black}18\times[75\div(19+6)]-25\\&\color{answerred}=18\times[75\div25]-25\\&\color{answerred}=18\times3-25\\&\color{answerred}=54-25\\&\color{answerred}=29\end{aligned}}
\CalcBlock{4}{\begin{aligned}[t]&\phantom{={}}\color{black}360\div[12\times(10-5)]+30\\&\color{answerred}=360\div[12\times5]+30\\&\color{answerred}=360\div60+30\\&\color{answerred}=6+30\\&\color{answerred}=36\end{aligned}}
\CalcBlock[简便解法（补充）]{5}{\begin{aligned}[t]&\phantom{={}}\color{black}12\times [36+(48-12)]\div 6\\&\color{answerred}=12\times [36+36]\div 6\\&\color{answerred}=12\times 72\div 6\\&\color{answerred}=(12\div 6)\times 72\\&\color{answerred}=2\times 72\\&\color{answerred}=144\end{aligned}}
\MethodBlock{按运算顺序计算（基础练习）}{本单元还要练习运算顺序；与上面的简便解法对照。}{\begin{aligned}[t]&\phantom{={}}12\times[36+(48-12)]\div6\\&=12\times[36+36]\div6\\&=12\times72\div6\\&=864\div6\\&=144\end{aligned}}



\CalcBlock{6}{\begin{aligned}[t]&\phantom{={}}\color{black}180\div[(76-28)\div4]+15\\&\color{answerred}=180\div[48\div4]+15\\&\color{answerred}=180\div12+15\\&\color{answerred}=15+15\\&\color{answerred}=30\end{aligned}}
\par\noindent

\QHeading{第2题　根据运算顺序，把下列每组算式分别合并成综合算式。}

\Problem{（1）\quad \(175-128=47\)\quad \(12\times4=48\)\quad \(47+48=95\)}
\par

\Answer{综合算式：$(175-128)+12\times4$}
\Answer{合并方法：用\((175-128)\)替换47，用\(12\times4\)替换48；小括号标明先算减法。}

\Problem{（2）\quad \(360\div60=6\)\quad \(15\times6=90\)\quad \(120-90=30\)}
\Answer{综合算式：$120-15\times(360\div60)$}
\Answer{合并方法：把90替换为\(15\times6\)，再把6替换为\(360\div60\)；给除法添上小括号，保持原来的运算顺序。}

\par\noindent

\QHeading{第3题　看图列式计算。}



\par
\Answer{\Formula{\begin{aligned}[t]&\phantom{={}}452-50\times4\\&=452-200\\&=252\text{（米）}\end{aligned}}\par 答：未知线段长252米。}




\par\noindent\makebox[\linewidth][c]{
\begin{tikzpicture}[x=1bp,y=-1bp,color=black,draw=black,text=black]
\begin{scope}\clip (297.0,650.0) rectangle (580.0,752.0);% 高清原题重排：PDF 9，原书 4。不含解答。
\hdhead{4}{四则混合运算——信息窗2}
\hdsection{1}{149}{脱式计算。}
\q{102}{193}{355}{13.5}{\(24\times[270\div(8+1)]-30\)}
\q{489}{193}{355}{13.5}{\([600-(120+80)]\div25\times4\)}
\q{102}{366}{355}{13.5}{\(18\times[75\div(19+6)]-25\)}
\q{489}{366}{355}{13.5}{\(360\div[12\times(10-5)]+30\)}
\q{102}{536}{355}{13.5}{\(12\times[36+(48-12)]\div6\)}
\q{489}{536}{355}{13.5}{\(180\div[(76-28)\div4]+15\)}
\hdsection{2}{718}{根据运算顺序，把下列每组算式分别合并成综合算式。}
\q{106}{765}{30}{13.5}{①}
\q{480}{762}{30}{13.5}{②}
\q{140}{765}{300}{13.5}{\(175-128=47\)}
\q{140}{807}{300}{13.5}{\(12\times4=48\)}
\q{140}{850}{300}{13.5}{\(47+48=95\)}
\q{512}{762}{300}{13.5}{\(360\div60=6\)}
\q{512}{804}{300}{13.5}{\(15\times6=90\)}
\q{512}{846}{300}{13.5}{\(120-90=30\)}
\q{139}{892}{115}{13.5}{综合算式：}
\q{512}{889}{115}{13.5}{综合算式：}
\begin{scope}[scale=.6666667,draw=black,line width=.45pt]
\draw (246,913)--(432,913) (619,910)--(808,910);
\end{scope}
\hdsection{3}{948}{看图列式计算。}
\q{106}{992}{30}{13.5}{①}
\q{450}{990}{30}{13.5}{②}
\q{139}{999}{68}{12}{50米}
\q{309}{991}{80}{12}{?米}
\q{232}{1054}{100}{12}{452米}
\q{496}{989}{126}{12}{150千克}
\q{481}{1013}{60}{12}{桃：}
\q{481}{1061}{60}{12}{梨：}
\q{584}{1038}{174}{12}{桃质量的5倍}
\q{756}{1009}{92}{12}{多100}
\q{756}{1035}{80}{12}{千克}
\q{625}{1096}{130}{12}{?千克}
\begin{scope}[scale=.6666667,draw=black,line width=.5pt]
\draw (142,1034)--(388,1034);
\foreach \x in {142,171,200,227,255,388}{\draw (\x,1027)--(\x,1034);}
\draw (143,1025)--(145,1021)--(153,1021)--(157,1018)--(162,1021)--(170,1021)--(171,1025);
\draw (255,1019)--(258,1015)--(318,1015)--(322,1012)--(326,1015)--(385,1015)--(388,1019);
\draw (143,1044)--(146,1048)--(261,1048)--(265,1051)--(269,1048)--(385,1048)--(388,1044);
\draw (518,1030)--(567,1030) (518,1023)--(518,1030) (567,1023)--(567,1030);
\draw (518,1017)--(521,1013)--(539,1013)--(543,1010)--(547,1013)--(564,1013)--(567,1017);
\draw (518,1078)--(797,1078);
\foreach \x in {518,567,615,664,713,763,797}{\draw (\x,1071)--(\x,1078);}
\draw (518,1067)--(521,1063)--(637,1063)--(641,1060)--(645,1063)--(760,1063)--(763,1067);
\draw (763,1067)--(766,1063)--(776,1063)--(780,1060)--(784,1063)--(794,1063)--(797,1067);
\draw (518,1084)--(521,1089)--(654,1089)--(658,1092)--(662,1089)--(794,1089)--(797,1084);
\end{scope}
\end{scope}\pgfresetboundingbox\path[use as bounding box] (297.0,650.0) rectangle (580.0,752.0);
\end{tikzpicture}}\par
\Answer{\Formula{\begin{aligned}[t]&\phantom{={}}150\times5+100\\&=750+100\\&=850\text{（千克）}\end{aligned}}\par 答：梨的质量是850千克。}

